\documentclass{article}
\usepackage{amsmath,amssymb}
\usepackage[margin=1in]{geometry}
\providecommand{\coloneqq}{\mathrel{:=}}
\providecommand{\llbracket}{[\![}
\providecommand{\rrbracket}{]\!]}

\begin{document}

\section*{A proposition as a function}

A proposition is a function $f^{\mathrm{prop}}$. We define it as a sequence
\[
  f^{\mathrm{prop}} = \langle f_1, f_2, \dots, f_n \rangle ,
\]
meaning it is evaluated step by step:
\begin{itemize}
  \item The domain of $f_1$ lies in the domain of $f^{\mathrm{prop}}$; apply the
        relevant terms to get $\operatorname{cod}(f_1)$.
  \item For each $k$, the domain of $f_k$ lies in
        $\operatorname{dom}(f^{\mathrm{prop}}) \cup \operatorname{cod}(f_1) \cup
        \dots \cup \operatorname{cod}(f_{k-1})$; apply the relevant terms to get
        $\operatorname{cod}(f_k)$.
  \item This continues until $\operatorname{cod}(f_n) = \operatorname{cod}(f^{\mathrm{prop}})$,
        the goal.
\end{itemize}

Equivalently:
\[
  \operatorname{dom}(f_k) \subseteq
    \operatorname{dom}(f^{\mathrm{prop}}) \cup
    \bigcup_{j<k} \operatorname{cod}(f_j),
  \qquad k = 1, \dots, n,
  \qquad
  \operatorname{cod}(f_n) = \operatorname{cod}(f^{\mathrm{prop}}).
\]


The union above is imprecise: it says $f_k$'s inputs lie \emph{somewhere} in
everything proved so far, but a step actually uses a \emph{specific} set of prior
facts. The two forms below name those dependencies exactly.

\section*{The standard way: a numbered derivation}

Fix the hypotheses $H = (h_1, \dots, h_m)$ (the ``given'' of the proposition).
A proof is a numbered list of facts. Each line states a fact $S_k$, the rule
$f_k$ that produces it, and the exact earlier lines / hypotheses it consumes:
\[
  \begin{array}{r@{\quad}l@{\qquad}l}
    s_1 : S_1 & \text{by } f_1 & \text{from } \{ \text{some of } H \} \\
    s_2 : S_2 & \text{by } f_2 & \text{from } \{ \text{some of } H,\, s_1 \} \\
    \vdots & \vdots & \vdots \\
    s_k : S_k & \text{by } f_k & \text{from } D_k \subseteq \{ H, s_1, \dots, s_{k-1} \} \\
    \vdots & \vdots & \vdots \\
    s_n : S_n & \text{by } f_n & \text{from } D_n, \qquad S_n = G .
  \end{array}
\]
Here $D_k$ is the \emph{precise} list of inputs to step $k$ (this is what your
$\mathtt{step}_k$ file records as its binders), the goal is $G$, and the proof is
complete because the last fact $S_n$ is $G$. This is the same sequence
$\langle f_1, \dots, f_n \rangle$ as above --- only the vague union
$\bigcup_{j<k}\operatorname{cod}(f_j)$ is replaced by the exact dependency set
$D_k$.

\section*{The same thing as a graph}

Equivalently, the proof is a directed acyclic graph. Nodes are the facts
$H, s_1, \dots, s_n$; a node $s_k$ has an incoming edge from each fact in its
dependency set $D_k$:
\[
  d \longrightarrow s_k \quad\text{for every } d \in D_k,
  \qquad s_n = G \text{ (the sink).}
\]
The numbered list above is just one topological ordering of this graph. Writing
$f^{\mathrm{prop}} = \langle f_1, \dots, f_n \rangle$ picks such an ordering and
then forgets the individual edges; the graph (or the list with its $D_k$'s) keeps
them.

\section*{The fully formal way: the proof term}

This is the precise object Lean stores. Nothing is a black box: every step names
the lemma it uses and the exact arguments it is applied to.

\subsection*{Notation, each symbol explained}

\begin{description}
  \item[$t : T$] ``$t$ has type $T$.'' Since a proposition \emph{is} a type,
    this reads ``$t$ is a proof of the proposition $T$.'' E.g.\ $h : A = B$ means
    $h$ is a proof that $A=B$.
  \item[$A \to B$] the type of \emph{functions} from $A$ to $B$. Read as a
    proposition it is the implication ``$A$ implies $B$'': a function that turns a
    proof of $A$ into a proof of $B$.
  \item[$\displaystyle\prod_{x : A} B(x)$] a \emph{dependent} function type
    (``for all $x : A$, $B(x)$''). Same as $A \to B$ except the output type
    $B(x)$ may mention the input $x$. This is how ``for all points $x$, \dots''
    is written. When $B$ does not mention $x$ it collapses to $A \to B$.
  \item[$f\, a$] \emph{application}: the function $f$ applied to the argument $a$.
    Just juxtaposition, no parentheses needed; $f\,a\,b$ means $(f\,a)\,b$, i.e.\
    $f$ applied to $a$ then to $b$. This is the ``what is applied to what.''
  \item[$\lambda x.\, e$] \emph{lambda abstraction}: the function that takes an
    input $x$ and returns $e$ (where $e$ may use $x$). The dot ``$.$'' just
    separates the input variable $x$ from the body $e$. So
    $\lambda x.\, x$ is the identity function, and $\lambda h.\, (\dots h \dots)$
    is ``given a proof $h$, build $\dots h \dots$.''
  \item[$\mathbf{let}\ x \coloneqq e_1\ \mathbf{in}\ e_2$] name the value $e_1$ as
    $x$, then use $x$ inside $e_2$. Pure abbreviation: it is exactly $e_2$ with
    $x$ standing for $e_1$. This is what lets us build one intermediate fact at a
    time.
  \item[$\langle a, b \rangle$] a \emph{pair} (tuple). As a proof, $\langle a,b\rangle$
    proves a conjunction ``$A \wedge B$'' by pairing a proof $a : A$ with a proof
    $b : B$. Used at the end to assemble several facts into the goal.
\end{description}

\subsection*{The proof term}

Let the proposition have hypotheses $h_1 : H_1,\ \dots,\ h_m : H_m$ and goal $G$.
Its type (the statement) is
\[
  f^{\mathrm{prop}} \;:\; \prod_{h_1 : H_1} \cdots \prod_{h_m : H_m} G ,
\]
and its proof (the term after ``$\coloneqq$'') is
\[
  f^{\mathrm{prop}} \;\coloneqq\;
  \lambda\, h_1\, \dots\, h_m.\;
  \begin{aligned}[t]
    &\mathbf{let}\ s_1 \coloneqq \ell_1(a_{1,1}, \dots, a_{1,r_1})\ \mathbf{in}\\
    &\mathbf{let}\ s_2 \coloneqq \ell_2(a_{2,1}, \dots, a_{2,r_2})\ \mathbf{in}\\
    &\quad\vdots\\
    &\mathbf{let}\ s_n \coloneqq \ell_n(a_{n,1}, \dots, a_{n,r_n})\ \mathbf{in}\\
    &s_n .
  \end{aligned}
\]
Reading it: each $\ell_k$ is a \emph{named} lemma (a cited proposition or an
axiom), and each argument $a_{k,j}$ is a \emph{specific} earlier fact --- one of
the hypotheses $h_1, \dots, h_m$ or an earlier step $s_1, \dots, s_{k-1}$. So
$s_k \coloneqq \ell_k(a_{k,1}, \dots, a_{k,r_k})$ says exactly ``apply lemma
$\ell_k$ to these particular inputs to obtain fact $s_k$.'' The vague $f_k$ from
before is now the concrete pair (name $\ell_k$, argument list $a_{k,\bullet}$),
and the dependency set is $D_k = \{ a_{k,1}, \dots, a_{k,r_k} \}$. The final line
returns $s_n$, whose type is the goal $G$.

\subsection*{A concrete step}

For instance a single step applying Euclid's Proposition~4 (SAS) might read
\[
  \mathbf{let}\ s_5 \coloneqq \mathtt{proposition\_4}\;(h_1)\,(s_2)\,(s_3)\ \mathbf{in}\ \dots
\]
meaning: the lemma is $\mathtt{proposition\_4}$; it is applied to hypothesis
$h_1$ and to the earlier facts $s_2$ and $s_3$; the result is named $s_5$. And
if the goal $G$ is a conjunction $G_1 \wedge G_2$ proved by facts $s_i$ and
$s_j$, the last line is $\langle s_i, s_j \rangle$ instead of a single $s_n$.

\section*{Naming the sequence of steps: the derivation}

The term $f^{\mathrm{prop}}$ is a \emph{single} object $\lambda \vec h.\,
(\mathbf{let}\ \dots\ \mathbf{in}\ s_n)$. Once assembled it is atomic: the
$\mathbf{let}$s reduce away, so there is \emph{no} projection $\pi_k(f^{\mathrm{prop}})$
that returns ``the $k$-th step.'' To speak of the ordered list of steps we keep a
separate, inspectable object and give it a name --- its \emph{derivation}.

\subsection*{New notation}

\begin{description}
  \item[$\sigma_k$] the $k$-th \emph{step}, i.e.\ the triple
    $\sigma_k = \langle \ell_k,\ D_k,\ S_k \rangle$: the lemma $\ell_k$ used, its
    dependency set $D_k \subseteq \{ \vec h, s_1, \dots, s_{k-1} \}$ (the exact
    inputs), and the fact $S_k$ it produces ($s_k : S_k$).
  \item[$\mathcal{D}$] the \emph{derivation}: the ordered tuple of all steps,
    \[
      \mathcal{D} = (\sigma_1, \sigma_2, \dots, \sigma_n).
    \]
    Unlike $f^{\mathrm{prop}}$, this \emph{is} indexable --- $\sigma_k$, $\ell_k$,
    $D_k$, $S_k$ all make sense.
  \item[$\llbracket \cdot \rrbracket$] the \emph{assembly} (elaboration) map,
    which collapses a derivation into a single proof term:
    \[
      f^{\mathrm{prop}} = \llbracket \mathcal{D} \rrbracket
      = \lambda\, \vec h.\;
        \mathbf{let}\ s_1 \coloneqq \ell_1(D_1)\ \mathbf{in}\ \cdots\
        \mathbf{let}\ s_n \coloneqq \ell_n(D_n)\ \mathbf{in}\ s_n .
    \]
\end{description}

\subsection*{The point}

$\mathcal{D}$ and $f^{\mathrm{prop}}$ are \emph{different objects}:
$\mathcal{D}$ is an ordered, inspectable list of steps, while
$f^{\mathrm{prop}} = \llbracket \mathcal{D} \rrbracket$ is the flat term it
elaborates to. You refer to individual steps through $\mathcal{D}$ (as $\sigma_k$,
$D_k$, $S_k$), never by trying to index into $f^{\mathrm{prop}}$. In our pipeline
each backing file $\mathtt{step}_k$ \emph{is} a $\sigma_k$ (its binders are $D_k$,
its theorem type is $S_k$), and wiring the proof is exactly computing
$\llbracket \mathcal{D} \rrbracket$.

\end{document}
